\documentclass[preview]{standalone}

\usepackage{amsmath}
\usepackage{amssymb}
\usepackage{stellar}
\usepackage{definitions}

\begin{document}

\id{free-modules-over-principal-ideal-domains}
\genpage

\section{Factorization length}

\begin{snippetdefinition}{pid-element-factorization-length}{Factorization length}
    Let \(A\) be a \principalidealdomain and let
    \(a\in A\difference\{0_A\}\). Write
    \[
        a=up_1^{n_1}\cdots p_r^{n_r},
        \qquad
        u\in A^\times,
    \]
    where the \(p_i\) are pairwise nonassociate
    \primeelement[prime elements]. The \emph{factorization length} of
    \(a\) is
    \[
        \ell(a)=n_1+\cdots+n_r.
    \]
\end{snippetdefinition}

\begin{snippetnote}{pid-element-factorization-length-properties}{Properties of factorization length}
    Unique factorization shows that \(\ell(a)\) is well-defined. If
    \(b\divides a\), then
    \[
        \ell(b)\leq\ell(a),
    \]
    because every prime exponent in \(b\) is at most the corresponding
    exponent in \(a\).
\end{snippetnote}

\section{Splitting along a homomorphism}

\begin{snippetlemma}{pid-module-homomorphism-splitting}{Splitting along a homomorphism to the PID}
    Let \(\varphi\colon M\fromto A\) be a \modulehomomorphism. Suppose
    that some \(v\in M\) satisfies
    \[
        \moduleimage(\varphi)=(\varphi(v)),
        \qquad
        \varphi(v)\neq0_A.
    \]
    Then
    \[
        M=Av\oplus\moduleker(\varphi).
    \]
\end{snippetlemma}

\begin{snippetproof}{pid-module-homomorphism-splitting-proof}{pid-module-homomorphism-splitting}{Splitting along a homomorphism to the PID}
    Put \(b=\varphi(v)\). If \(av\in\moduleker(\varphi)\), then
    \[
        0_A=\varphi(av)=ab.
    \]
    Since \(A\) is an \integraldomain and \(b\neq0_A\), one has \(a=0_A\).
    Thus \(Av\intersection\moduleker(\varphi)=\{0_M\}\).

    For any \(w\in M\), the equality
    \(\moduleimage(\varphi)=(b)\) gives \(\varphi(w)=cb\) for some
    \(c\in A\). Hence
    \[
        \varphi(w-cv)=\varphi(w)-c\varphi(v)=0_A,
    \]
    and
    \[
        w=cv+(w-cv)\in Av+\moduleker(\varphi).
    \]
    The sum is therefore direct and equals \(M\).
\end{snippetproof}

\section{Submodules of free modules}

\begin{snippettheorem}{pid-submodule-of-finite-free-module}{Submodules of finite free modules over a PID}
    Let \(M\) be a finite \freemodule over the
    \principalidealdomain \(A\), with
    \[
        \freemodulerank(M)=r.
    \]
    Every \submodule \(V\submoduleleq M\) is free and
    \[
        \freemodulerank(V)\leq r.
    \]
\end{snippettheorem}

\begin{snippetproof}{pid-submodule-of-finite-free-module-proof}{pid-submodule-of-finite-free-module}{Submodules of finite free modules over a PID}
    Proceed by induction on \(r\). The case \(r=0\) is immediate. If
    \(r=1\), identify \(M\) with \(A\). Then \(V\) corresponds to an
    \ideal \(I=(c)\). If \(c=0_A\), it is free of rank \(0\); otherwise,
    multiplication by \(c\) defines an isomorphism \(A\to(c)\), because
    \(A\) is a domain. Thus \(V\) is free of rank \(1\).

    Suppose \(r>1\) and identify
    \[
        M=A^{r-1}\oplus A.
    \]
    Put \(W=A^{r-1}\oplus\{0_A\}\) and let
    \(\pi\colon M\fromto A\) be projection onto the last coordinate.
    By induction, \(V\intersection W\) is free of rank at most \(r-1\).
    If \(V\subseteq W\), the result follows.

    Otherwise \(\pi|_V\) has nonzero image. Since \(A\) is a PID,
    \[
        \moduleimage(\pi|_V)=(c)
    \]
    for some \(c\neq0_A\). Choose \(v\in V\) with \(\pi(v)=c\). The
    splitting lemma applied to \(\pi|_V\) gives
    \[
        V=Av\oplus\moduleker(\pi|_V)
        =
        Av\oplus(V\intersection W).
    \]
    Hence \(V\) is free of rank at most \(r\).
\end{snippetproof}

\section{Primitive vectors and coefficient gcds}

\begin{snippetlemma}{free-module-coefficient-gcd-homomorphisms}{Coefficient gcds and homomorphisms to the PID}
    Let \(M\) be a finite \freemodule with basis
    \(\mathcal X=\{v_1,\ldots,v_r\}\), and write
    \[
        w=a_1v_1+\cdots+a_rv_r.
    \]
    Let \(b\) be a greatest common divisor of the coefficients, so
    \[
        (b)=(a_1,\ldots,a_r).
    \]
    Then:
    \begin{enumerate}
        \item \(b\divides\varphi(w)\) for every
        \modulehomomorphism \(\varphi\colon M\fromto A\);
        \item there exists a \modulehomomorphism
        \(\psi\colon M\fromto A\) such that \(\psi(w)=b\).
    \end{enumerate}
\end{snippetlemma}

\begin{snippetproof}{free-module-coefficient-gcd-homomorphisms-proof}{free-module-coefficient-gcd-homomorphisms}{Coefficient gcds and homomorphisms to the PID}
    For every \(\varphi\colon M\fromto A\),
    \[
        \varphi(w)=a_1\varphi(v_1)+\cdots+a_r\varphi(v_r).
    \]
    Since \(b\) divides every \(a_i\), it divides \(\varphi(w)\).

    The identity \((b)=(a_1,\ldots,a_r)\) gives coefficients
    \(c_1,\ldots,c_r\in A\) such that
    \[
        b=a_1c_1+\cdots+a_rc_r.
    \]
    By the universal property of the free module, there is a unique
    \modulehomomorphism \(\psi\colon M\fromto A\) with
    \(\psi(v_i)=c_i\). It satisfies \(\psi(w)=b\).
\end{snippetproof}

\begin{snippetnote}{free-module-coefficient-gcd-associates}{The coefficient gcd is defined up to associates}
    The greatest common divisor in the previous lemma is understood up
    to multiplication by a unit. Equivalently, it is any generator of
    the \ideal \((a_1,\ldots,a_r)\).
\end{snippetnote}

\section{Compatible bases}

\begin{snippettheorem}{pid-submodule-compatible-bases}{Diagonal form for a submodule of a free module}
    Let \(M\) be a free \(A\)-module of rank \(r\), and let
    \(V\submoduleleq M\). There exist a basis
    \[
        \mathcal X=\{x_1,\ldots,x_r\}
    \]
    of \(M\) and elements \(c_1,\ldots,c_r\in A\) such that
    \[
        V=\langle c_1x_1,\ldots,c_rx_r\rangle.
    \]
    After discarding the zero generators, the remaining vectors form a
    basis of \(V\).
\end{snippettheorem}

\begin{snippetproof}{pid-submodule-compatible-bases-proof}{pid-submodule-compatible-bases}{Diagonal form for a submodule of a free module}
    Proceed by induction on \(r\). The cases \(r=0\) and \(V=\{0_M\}\)
    are immediate. Fix an initial basis
    \(\mathcal Y=\{y_1,\ldots,y_r\}\). For each nonzero
    \[
        z=a_1y_1+\cdots+a_ry_r\in V,
    \]
    let \(d(z)\) be a gcd of \(a_1,\ldots,a_r\). Choose \(w\neq0_M\)
    in \(V\) so that \(\ell(d(w))\) is minimal. Write
    \[
        w=a_1y_1+\cdots+a_ry_r,
        \qquad
        a_i=da_i',
        \qquad
        d=d(w),
    \]
    and put \(x_1=\sum_i a_i'y_i\). Then \(w=dx_1\) and
    \((a_1',\ldots,a_r')=A\). The gcd lemma gives a
    \modulehomomorphism \(\psi\colon M\fromto A\) such that
    \(\psi(x_1)=1_A\). The splitting lemma yields
    \[
        M=Ax_1\oplus\moduleker(\psi).
    \]

    Now \(\moduleimage(\psi|_V)=(c)\) for some \(c\in A\). Since
    \(\psi(w)=d\), one has \(c\divides d\). Choose \(u\in V\) with
    \(\psi(u)=c\). The gcd lemma gives \(d(u)\divides c\), while the
    minimal choice of \(w\) gives
    \[
        \ell(d)\leq\ell(d(u)).
    \]
    Consequently,
    \[
        \ell(d)
        \leq\ell(d(u))
        \leq\ell(c)
        \leq\ell(d).
    \]
    Thus \(c\) and \(d\) are associates, and
    \[
        \moduleimage(\psi|_V)=(d)=(\psi(w)).
    \]
    Applying the splitting lemma to \(\psi|_V\) gives
    \[
        V=Aw\oplus\moduleker(\psi|_V)
        =
        A(dx_1)\oplus(V\intersection\moduleker(\psi)).
    \]

    The kernel of \(\psi\) is free of rank \(r-1\). By induction, it has
    a basis \(\{x_2,\ldots,x_r\}\) and there are
    \(c_2,\ldots,c_r\in A\) such that
    \[
        V\intersection\moduleker(\psi)
        =
        \langle c_2x_2,\ldots,c_rx_r\rangle.
    \]
    Taking \(c_1=d\), the basis
    \(\{x_1,\ldots,x_r\}\) of \(M\) has the required form.
\end{snippetproof}

\begin{snippetcorollary}{full-rank-integer-submodule-index}{Index of a full-rank integer submodule}
    Let \(M\) be a finite free \(\integers\)-module, and let
    \(V\submoduleleq M\) satisfy
    \[
        \freemodulerank(V)=\freemodulerank(M)=r.
    \]
    If \(\mathcal X=\{x_1,\ldots,x_r\}\) is a basis of \(M\),
    \(\mathcal Y=\{y_1,\ldots,y_r\}\) is a basis of \(V\), and
    \(C\in\operatorname{Mat}_r(\integers)\) expresses the vectors of
    \(\mathcal Y\) in the basis \(\mathcal X\), then
    \[
        [M:V]=|\det C|.
    \]
\end{snippetcorollary}

\begin{snippetproof}{full-rank-integer-submodule-index-proof}{full-rank-integer-submodule-index}{Index of a full-rank integer submodule}
    Compatible bases give a basis \(\{x_1,\ldots,x_r\}\) of \(M\) and
    nonzero integers \(c_1,\ldots,c_r\) such that
    \[
        V=\langle c_1x_1,\ldots,c_rx_r\rangle.
    \]
    Hence
    \[
        M/V
        \moduleisomorphic
        \integers/(c_1)\oplus\cdots\oplus\integers/(c_r)
    \]
    and
    \[
        [M:V]=|M/V|=|c_1\cdots c_r|.
    \]
    In these bases, the inclusion matrix is
    \(\operatorname{diag}(c_1,\ldots,c_r)\). Arbitrary changes of bases
    multiply it on the left and right by unimodular matrices, whose
    determinants are \(\pm1\). Thus the absolute value of the
    determinant remains \(|c_1\cdots c_r|\).
\end{snippetproof}

\end{document}
